% Page 025

\seite{25}

Am Schlüsse des Schuljahres wurden drei Kinder\ob
der Schule entlassen. Zwei Kinder.

\pend
\jahresueberschrift{Das Jahr 1897/98}
\pstart

Ostern zu Ostern 1897 5 Schulanfänger\ob
aufgenommen. Die Schülerzahl beträgt jetzt\ob
61.\ Sämtliche Kinder sind katholisch und wohl\ohb
fährig. Auch den letzten Herbst werden die\ob
Schulräumen zueng; nöthig die Frage über einen\ob
Neubau immer mehr in den Vordergrund tritt.

\pend
\abschnitt{Pfarrwechsel}
\pstart

Mit Beginn des Monats Mai wurde Herr Pfarrer\ob
Michael Dunckel, der auch zugleich Lokalschul\ohb
inspektor war, nach Karwesheim\unsicher{} versetzt. Die Pfarrei hat den\ob
Herrn mit ungern scheiden, da er sich als uneigennütziger\ob
und unserer Kirchenfreund bei Gegen\ohb
über im Kürze gewonnen. Sein Nachfolger\ob
wurde der bisherige Kaplan in W.\ Paulin Fritz\ob
zur Ettges, der am 13.\ Mai 1897 feier\ohb
lich in sein Amt eingeführt wurde.

Mit Beginn des Sommersemesters hat ein jeder sich\ob
hiesige angehende Bürger ins Plüstelavin\unsicher{} zu hier\ob
um sich den erfahrenen Kursus zu nehmen.\ob
Diese Bestrebung ist es nicht doch Winterabendschule besitzt\ob
dem Lehrer das Gymnasium in Prüm und über.\ob
Für junge Leute, die augenblicklich auf Privatstücke\ob
sich genießen, sind auch entschlossen diesen\ob
Kursus zu nehmen seiner dieser aufschrift befremd.
\pend
